\subsection{箭图中的道路与圈}

定义 2. —— 设 C 为箭图。C 中的道路是一个序列 \(c = (a_{0}, f_{1}, a_{1}, \ldots, a_{n-1}, f_{n}, a_{n})\) ，其中 \(n\) 是整数 \(\geqslant 0\) ，\(a_{0}, \ldots, a_{n}\) 是 C 的顶点，且对 \(1 \leqslant i \leqslant n\) ，\(f_{i}\) 是 C 中连接 \(a_{i-1}\) 与 \(a_{i}\) 的箭。我们说道路 c 的长度为 \(n\) 。

设 \(c = (a_{0}, f_{1}, a_{1}, \ldots, a_{n-1}, f_{n}, a_{n})\) 为 C 中长度为 n 的道路。顶点 \(a_{0}\) 称为 c 的起点或源；顶点 \(a_{n}\) 称为 c 的终点或靶；也说 c 连接顶点 \(a_{0}\) 与顶点 \(a_{n}\) 。顶点 \(a_{i}\) （\(0 \leqslant i \leqslant n\)）称为指标 i 的顶点；箭 \(f_{i}\) （\(1 \leqslant i \leqslant n\)）称为第 i 个箭，或指标 i 的箭。长度 \(\geqslant 1\) 的道路由其箭的序列确定。

起点与终点相同的道路称为圈。长度为 0 的道路是圈，称为常值圈。

我们说道路

\[
c = (a _ {0}, f _ {1}, a _ {1}, \ldots , a _ {n - 1}, f _ {n}, a _ {n}) \text{且} c ^ {\prime} = (a _ {0} ^ {\prime}, f _ {1} ^ {\prime}, a _ {1} ^ {\prime}, \ldots , a _ {m - 1} ^ {\prime}, f _ {m} ^ {\prime}, a _ {m} ^ {\prime})
\]

在 C 中是可并列的，如果 c 的终点 \(a_{n}\) 是起点 \(a_{0}^{\prime}\) （\(c^{\prime}\) 的）。这时，序列 \((a_{0}, f_{1}, a_{1}, \ldots, a_{n-1}, f_{n}, a_{n}, f_{1}^{\prime}, a_{1}^{\prime}, \ldots, f_{m}^{\prime}, a_{m}^{\prime})\) 是 C 中的道路，记作 \(c * c^{\prime}\) ，称为 c 与 \(c^{\prime}\) 的并列道路。它连接 c 的起点与 \(c^{\prime}\) 的终点；其长度是道路 c 与 \(c^{\prime}\) 的长度之和。

